\documentclass[a4paper,11pt]{article}

\usepackage[top=2cm,bottom=2cm,left=2.5cm,right=2.5cm]{geometry}
\usepackage{fontspec}

\setmainfont{TeX Gyre Pagella}

% Symbola dekkjer heile Unicode Geometric Shapes Extended + miscellaneous
\newfontfamily{\symbolfont}{Symbola}[Scale=0.95]
\newcommand{\g}[1]{{\symbolfont #1}}

\usepackage{unicode-math}
\setmathfont{TeX Gyre Pagella Math}

\usepackage{microtype}
\usepackage{array}
\usepackage{titlesec}

\titleformat{\section}{\normalfont\scshape\large}{\S\thesection}{1em}{}
\titleformat{\subsection}{\normalfont\scshape\normalsize}{\thesubsection}{1em}{}
\titlespacing*{\section}{0pt}{1.5em}{0.7em}
\titlespacing*{\subsection}{0pt}{1em}{0.4em}

\pagestyle{plain}
\setlength{\parindent}{0em}
\setlength{\parskip}{0.35em}

% ================================================================
% UNICODE-GLYFAR
% Dei blir alle skrivne med \g{…} slik at dei bruker DejaVu Sans
% ================================================================
\newcommand{\gAgent}{\g{▲}}
\newcommand{\gDesignar}{\g{(▲)}}
\newcommand{\gForm}{\g{●}}
\newcommand{\gGrammatikk}{\g{⬢}}
\newcommand{\gFormrom}{\g{□}}
\newcommand{\gKlasse}{\g{\{●\}}}
\newcommand{\gLyskjegle}{\g{◇}}
\newcommand{\gKonfig}{\g{⁂}}
\newcommand{\gObjekt}{\g{•}}
\newcommand{\gOmhug}{\g{☉}}
\newcommand{\gPartisjon}{\g{◐}}
\newcommand{\gTrykk}{\g{↘}}
\newcommand{\gStilart}{\g{(•••)}}
\newcommand{\gSubstrat}{\g{▽}}
\newcommand{\gLandskap}{\g{∿}}
\newcommand{\gTransform}{\g{⬟}}
\newcommand{\gSirkel}{\g{◯}}

\begin{document}

\thispagestyle{empty}

\begin{center}
{\Large\scshape Formskrift}\\[0.2em]
{\itshape Formalisert alfabet, kompositorreglar og bevis}
\end{center}

\vspace{0.5em}

\begin{center}
\rule{\textwidth}{0.4pt}
\end{center}

% ================================================================
\section{Den fulle symbolmengda}

Formskrifta består av tre disjunkte tabellar: \textbf{formskrift-glyfar} (dei sekstan omgrepa), \textbf{operatorar} (relasjonar mellom glyfar), og \textbf{logikk} (kvantifikasjon og slutning). Saman utgjer dei eit komplett system der kvar proposisjon i Formlæra let seg uttrykke.

\subsection{Formskrift-glyfar}

\begin{center}
\begin{tabular}{@{}clc@{\hspace{1.5em}}cl@{}}
\textbf{Glyf} & \textbf{Omgrep} & & \textbf{Glyf} & \textbf{Omgrep}\\
\hline
\gAgent       & agent              & & \gOmhug       & omhug\\
\gDesignar    & designar           & & \gPartisjon   & partisjon\\
\gForm        & form               & & \gTrykk       & seleksjonstrykk\\
\gGrammatikk  & formgrammatikk     & & \gStilart     & stilart\\
\gFormrom     & formrom            & & \gSubstrat    & substrat\\
\gKlasse      & klasse             & & \gLandskap    & tilpassingslandskap\\
\gLyskjegle   & kognitiv lyskjegle & & \gTransform   & transformasjon\\
\gKonfig      & konfigurasjon      & & \gObjekt      & objekt\\
\end{tabular}
\end{center}

\subsection{Operatorar}

\begin{center}
\begin{tabular}{@{}cl@{\hspace{1.5em}}cl@{\hspace{1.5em}}cl@{}}
$\subset$         & delmengd            & $\nabla$  & gradient    & $\mid$      & vilkår\\
$\in$             & element             & $\rightarrow$ & implikasjon & $\cap$  & skjering\\
$\leftrightarrow$ & biimplikasjon       & $\triangleright$ & applikasjon & $|\cdot|$ & storleik\\
$=$               & likskap             & $\sum$    & sum         & & \\
\end{tabular}
\end{center}

\subsection{Logikk}

\begin{center}
\begin{tabular}{@{}cl@{\hspace{1em}}cl@{\hspace{1em}}cl@{}}
$\exists$       & eksistenskvantor  & $\forall$   & allkvantor   & $\therefore$ & slutning\\
$\nexists$      & negert eksistens  & $\because$  & premiss      & $\equiv$     & identitet\\
$\vee$          & disjunksjon       & $\wedge$    & konjunksjon  & $\neg$       & negasjon\\
$\varnothing$   & tom mengd         & $\cup$      & union        & & \\
\end{tabular}
\end{center}

% ================================================================
\section{Kompositorreglar}

Glyfane er ikkje alle primitive. Somme er samansette av enklare glyfar gjennom definerte kompositorreglar. Dette gjer formskrifta til ein \emph{algebra}: strukturar kan dekomponerast til dei enklare einingane dei er bygde av.

\subsection{Geometriske primitivar}

Kvar glyf høyrer til ein geometrisk familie som ber sitt eige semantiske innhald:

\begin{center}
\begin{tabular}{@{}cll@{}}
\textbf{Familie} & \textbf{Glyfar} & \textbf{Semantikk}\\
\hline
Trekant (3-kant)  & \gAgent\ / \gSubstrat                & agens / substrat (dualpar)\\
Firkant (4-kant)  & \gFormrom                             & formrom\\
Femkant (5-kant)  & \gTransform                           & transformasjon, dynamikk\\
Sekskant (6-kant) & \gGrammatikk                          & formgrammatikk, regelsystem\\
Diamant (4-kant)  & \gLyskjegle                           & representasjonsrom\\
Sirkel            & \gSirkel\ / \gForm\ / \gPartisjon\ / \gOmhug & rom, form, partisjon, omhug\\
\end{tabular}
\end{center}

\subsection{Samansettingsreglar}

\begin{center}
\begin{tabular}{@{}rl@{\hspace{1em}}l@{}}
\textbf{R1.} & \gAgent\ $+$ \gSirkel\ $=$ \gDesignar         & agent $+$ sosial kontekst $=$ designar\\
\textbf{R2.} & \gSirkel\ $+$ \gObjekt\ $=$ \gOmhug           & open $+$ objekt-aksar $=$ omhug\\
\textbf{R3.} & \gForm\ $+$ \gSirkel\ $=$ \gPartisjon         & kjent $+$ tenkeleg $=$ partisjon\\
\textbf{R4.} & \gObjekt\ $+$ \gObjekt\ $+$ \gObjekt\ $=$ \gKonfig & tre objekt bundne $=$ konfigurasjon\\
\textbf{R5.} & $\{\,\cdot\,\}$ $+$ \gForm\ $=$ \gKlasse       & avgrensing $+$ typisk form $=$ klasse\\
\textbf{R6.} & $(\,\cdot\,)$ $+$ \gObjekt\gObjekt\gObjekt\ $=$ \gStilart & avgrensing $+$ aggregat $=$ stilart\\
\textbf{R7.} & $\nabla$\gLandskap\ $=$ \gTrykk              & gradient av landskap $=$ trykk\\
\textbf{R8.} & \gAgent\ $+$ \gSubstrat\ $\therefore$ \gLyskjegle & agent $+$ substrat gjev lyskjegle\\
\end{tabular}
\end{center}

Reglane har dobbelt funksjon. \emph{Didaktisk} gjer dei alfabetet lettare å læra: lesaren treng ikkje memorere sekstan uavhengige glyfar, berre åtte kombinasjonsreglar. \emph{Semantisk} viser dei kva omgrepa er bygde av. Designaren er ikkje eit nytt vesen; han er ein agent i sosial kontekst. Omhugen er ikkje mystisk; han er ein open kapasitet fylt av objekt-aksar.

\newpage

% ================================================================
\section{Bevis av sentrale teorem}

No som systemet er formalisert, kan vi \emph{rekne} med han. Vi bevisar sytten teorem frå traktaten ved reint formelle skritt. Kvart bevis brukar berre definisjonar, kompositorreglar og logikk-reglar frå dei tre tabellane.

\subsection{T1. Designaren er ein spesiell agent}

\emph{Påstand.} Alle designarar er agentar; ikkje alle agentar er designarar.

\begin{center}
\begin{tabular}{@{}rl@{\hspace{1em}}l@{}}
1. & \gDesignar\ $\equiv$ \gAgent\ $+$ \gSirkel                          & (R1, definisjon)\\
2. & $\therefore$ \gDesignar\ $\rightarrow$ \gAgent                      & (projeksjon)\\
3. & \gAgent\ $\not\rightarrow$ \gDesignar                               & (mangel av sosial \gSirkel)\\
4. & $\therefore$ \gDesignar\ $\subset$ \gAgent,\ \gDesignar\ $\neq$ \gAgent & (proposisjon 6.021)\\
\end{tabular}
\end{center}

\subsection{T2. Kjerne-likninga}

\emph{Påstand.} Forma er grammatikken sin respons på gradienten i landskapet, projisert på agenten si lyskjegle.

\begin{center}
\begin{tabular}{@{}rl@{\hspace{1em}}l@{}}
1. & \gForm\ $\in$ \gFormrom                                              & (proposisjon 1.3)\\
2. & \gGrammatikk\ $\triangleright$ \gFormrom                              & (proposisjon 5.3)\\
3. & \gTrykk\ $=$ $\nabla$\gLandskap                                      & (R7)\\
4. & \gLyskjegle\ $\subset$ \gFormrom                                     & (proposisjon 5.2)\\
5. & $\therefore$ \gForm\ $=$ \gGrammatikk$(\nabla$\gLandskap$)$ $\cap$ \gLyskjegle & (proposisjon 5.7)\\
\end{tabular}
\end{center}

\subsection{T3. Det etiske rommet}

\emph{Påstand.} Det etiske rommet i design er storleiken av lyskjegla minus omhugen.

\begin{center}
\begin{tabular}{@{}rl@{\hspace{1em}}l@{}}
1. & \gOmhug\ $\subset$ \gLyskjegle                                       & (proposisjon 5.6)\\
2. & $|$\gOmhug$|$ $\leq$ $|$\gLyskjegle$|$                              & (storleiksmonotoni)\\
3. & $|$\gLyskjegle$|$ $-$ $|$\gOmhug$|$ $\geq 0$                        & (aritmetikk)\\
4. & $\therefore$ \emph{etisk rom} $\equiv$ $|$\gLyskjegle$|$ $-$ $|$\gOmhug$|$ & (proposisjon 7.23)\\
\end{tabular}
\end{center}

\subsection{T4. Stilen har ikkje kausal kraft}

\emph{Påstand.} Stilarten er mønsterminne, ikkje ei kraft.

\begin{center}
\begin{tabular}{@{}rl@{\hspace{1em}}l@{}}
1. & \gStilart\ $=$ $\{$\gForm\ $|$ $d <$ $\varepsilon\}$                & (proposisjon 3.3)\\
2. & \gStilart\ $\equiv$ aggregat av former, ikkje kraft                  & (type-analyse)\\
3. & \gTrykk\ $\in$ krefter; \gStilart\ $\in$ aggregat                   & (disjunkte typar)\\
4. & $\therefore$ \gStilart\ $\not\rightarrow$ \gForm                    & (proposisjon 6.3)\\
\end{tabular}
\end{center}

\subsection{T5. Formrommet er epistemisk partisjonert}

\emph{Påstand.} Formrommet deler seg i kjent, tenkeleg og forbode.

\begin{center}
\begin{tabular}{@{}rl@{\hspace{1em}}l@{}}
1. & \gFormrom\ $=$ \gForm\ $\cup$ \gSirkel\ $\cup\ \varnothing_f$        & (R3, utvida)\\
2. & \gForm\ $\cap$ \gSirkel\ $=$ $\varnothing$                          & (disjunksjon)\\
3. & \gForm\ $\cap\ \varnothing_f$ $=$ $\varnothing$, \gSirkel\ $\cap\ \varnothing_f$ $=$ $\varnothing$ & (disjunksjon)\\
4. & $\therefore$ \gFormrom\ $=$ \gPartisjon\ $\cup\ \varnothing_f$       & (proposisjon 1.44)\\
\end{tabular}
\end{center}

\subsection{T6. Reduksjonsteoremet}

\emph{Påstand.} Når landskapet er flatt, reduserer kjerne-likninga til rein grammatikk.

\begin{center}
\begin{tabular}{@{}rl@{\hspace{1em}}l@{}}
1. & \gForm\ $=$ \gGrammatikk$(\nabla$\gLandskap$)$ $\cap$ \gLyskjegle   & (T2)\\
2. & $\nabla$\gLandskap\ $\equiv 0$ (føresetnad)                         & (vilkår)\\
3. & \gGrammatikk$(0)$ $=$ $L($\gGrammatikk$)$                           & (heile språket)\\
4. & $\therefore$ \gForm\ $=$ $L($\gGrammatikk$)$ $\cap$ \gLyskjegle     & (proposisjon 5.72)\\
\end{tabular}
\end{center}

\subsection{T7. Substratet som null-ordens agent}

\emph{Påstand.} Substratet fungerer som ein agent utan intensjonalitet.

\begin{center}
\begin{tabular}{@{}rl@{\hspace{1em}}l@{}}
1. & \gAgent\ $=$ $(g, d, \delta)$                                         & (proposisjon 5.11)\\
2. & \gSubstrat\ har $d$ og $\delta$, men $g = \varnothing$                & (observasjon)\\
3. & $\therefore$ \gSubstrat\ $\subset$ \gAgent, $g$-komponent $=\ \varnothing$ & (proposisjon 5.4)\\
\end{tabular}
\end{center}

\subsection{T8. Lyskjegla som skjering}

\emph{Påstand.} Poly-agentisk form ligg i skjeringa av alle lyskjegler.

\begin{center}
\begin{tabular}{@{}rl@{\hspace{1em}}l@{}}
1. & $\forall i:$ \gLyskjegle$_i$ $\subset$ \gFormrom                    & (proposisjon 5.2)\\
2. & \gForm\ må vere representert av kvar \gAgent$_i$                     & (proposisjon 6.1)\\
3. & $\therefore$ \gForm\ $\in\ \cap_i$ \gLyskjegle$_i$                   & (skjering)\\
\end{tabular}
\end{center}

\subsection{T9. Formrommet ekspanderer}

\emph{Påstand.} Kvar ny realisering utvider det kjende.

\begin{center}
\begin{tabular}{@{}rl@{\hspace{1em}}l@{}}
1. & \gForm$_{\text{ny}}$ realisert ved tid $t$                            & (hending)\\
2. & $K(t) = K(t{-}1) \cup \{$\gForm$_{\text{ny}}\}$                       & (akkumulering)\\
3. & $|K(t)| > |K(t{-}1)|$                                                 & (aritmetikk)\\
4. & $\therefore$ \gPartisjon\ er dynamisk                                 & (proposisjon 4.4)\\
\end{tabular}
\end{center}

\subsection{T10. Svikt-signaturen}

\emph{Påstand.} Ein form designa under trong omhug sviktar langs ikkje-representerte aksar.

\begin{center}
\begin{tabular}{@{}rl@{\hspace{1em}}l@{}}
1. & $\exists$ aks $a:$ $a \in$ \gLyskjegle, $a \notin$ \gOmhug           & (føresetnad)\\
2. & \gForm\ navigert kun etter aksar i \gOmhug                            & (T3)\\
3. & $\therefore$ \gForm\ optimaliserer ikkje for $a$                      & (definisjon)\\
4. & $\therefore$ \gForm\ sviktar langs $a$                                & (proposisjon 7.1)\\
\end{tabular}
\end{center}

\subsection{T11. Dynamikkens fundamentalteorem}

\emph{Påstand.} Landskapet er alltid i rørsle.

\begin{center}
\begin{tabular}{@{}rl@{\hspace{1em}}l@{}}
1. & \gLandskap\ $= \sum_i$ \gTrykk$_i$                                    & (proposisjon 3.1)\\
2. & $\exists i:$ $\partial$\gTrykk$_i$ $/\partial t \neq 0$               & (postulat 4.1)\\
3. & $\therefore$ $\partial$\gLandskap\ $/\partial t \neq 0$               & (distributiv)\\
4. & $\therefore$ \gTransform$($\gLandskap$)$ er ikkje-triviell            & (dynamisk system)\\
\end{tabular}
\end{center}

\subsection{T12. Blindskapens pris}

\emph{Påstand.} Ein-dimensjonal optimalisering produserer svikt langs alle andre dimensjonar.

\begin{center}
\begin{tabular}{@{}rl@{\hspace{1em}}l@{}}
1. & $\exists$ trykk $t^*:$ $t^*$ dominerer alle andre \gTrykk$_i$        & (føresetnad 4.5)\\
2. & \gOmhug\ kollapsar til $\{t^*\}$                                     & (konsekvens)\\
3. & $\forall$ aks $a \neq t^*:$ $a \notin$ \gOmhug                       & (T10-premiss)\\
4. & $\therefore$ \gForm\ sviktar langs $\forall a \neq t^*$              & (proposisjon 4.511)\\
\end{tabular}
\end{center}

\subsection{T13. Klasse og affordanse}

\emph{Påstand.} Ein klasse er definert av delt affordansestruktur, ikkje av ytre form.

\begin{center}
\begin{tabular}{@{}rl@{\hspace{1em}}l@{}}
1. & \gKlasse\ $=$ $\{$\gForm$_i$ $|$ $\mathrm{aff}(i) = \mathrm{aff}_0\}$ & (R5, presisert)\\
2. & \gKlasse\ $\neq$ $\{$\gForm$_i$ $|$ visuell likskap$\}$               & (ikkje fenomenologisk)\\
3. & $\therefore$ to \gForm\ i same \gKlasse\ kan sjå ulike ut            & (proposisjon 2.14)\\
\end{tabular}
\end{center}

\subsection{T14. Tripleteforeldreskap}

\emph{Påstand.} Kjerne-likninga har tre strukturelle foreldre.

\begin{center}
\begin{tabular}{@{}rl@{\hspace{1em}}l@{}}
1. & \gForm\ $=$ \gGrammatikk$(\nabla$\gLandskap$)$ $\cap$ \gLyskjegle   & (T2)\\
2. & Komponentar: \gGrammatikk, \gLandskap, \gLyskjegle                   & (dekomposisjon)\\
3. & Kvar komponent har eiga seksjon i traktaten                          & (observasjon)\\
4. & $\therefore$ \gForm\ er kausal konfluens av tre strukturar           & (§5.7 i traktat)\\
\end{tabular}
\end{center}

\subsection{T15. Agens er substrat-uavhengig}

\emph{Påstand.} Definisjonen av agent krev ingen bestemt fysisk realisering.

\begin{center}
\begin{tabular}{@{}rl@{\hspace{1em}}l@{}}
1. & \gAgent\ $=$ $(g, d, \delta)$                                         & (proposisjon 5.11)\\
2. & $(g, d, \delta)$ refererer til \emph{struktur}, ikkje \emph{materie} & (type-analyse)\\
3. & $\forall$ \gSubstrat\ som realiserer $(g, d, \delta):$ \gSubstrat\ $\rightarrow$ \gAgent & (R8)\\
4. & $\therefore$ \gAgent\ er substrat-uavhengig                           & (proposisjon 5.111)\\
\end{tabular}
\end{center}

\subsection{T16. Konvergens av tradisjonar}

\emph{Påstand.} Shape grammar og CK-teori er grensetilfelle av same likning.

\begin{center}
\begin{tabular}{@{}rl@{\hspace{1em}}l@{}}
1. & \gForm\ $=$ \gGrammatikk$(\nabla$\gLandskap$)$ $\cap$ \gLyskjegle    & (T2)\\
2. & Ved $\nabla$\gLandskap\ $\equiv 0:$ rein grammatikk (Stiny-Gips)     & (T6)\\
3. & Ved \gGrammatikk\ $\equiv \mathrm{id}:$ partisjon-operasjonar (Hatchuel-Weil) & (dualt grensetilfelle)\\
4. & $\therefore$ begge tradisjonar $\subset$ Formlæra                     & (proposisjon 5.722)\\
\end{tabular}
\end{center}

\subsection{T17. Stilart som makroskala-attraktor}

\emph{Påstand.} Stilen kanaliserer lyskjegla til observatøren.

\begin{center}
\begin{tabular}{@{}rl@{\hspace{1em}}l@{}}
1. & \gStilart\ $=$ formsverm kring tyngdepunkt                            & (T4)\\
2. & Observatøren sin \gLyskjegle\ lærer av realiserte \gForm              & (proposisjon 5.5)\\
3. & \gStilart\ utgjer konsentrert læringsgrunnlag                         & (observasjon)\\
4. & $\therefore$ \gStilart\ smalar \gLyskjegle\ til observatøren          & (proposisjon 6.3)\\
\end{tabular}
\end{center}

% ================================================================
\section{Refleksjon}

Sytten bevis, alle bygde av dei same operatorane og logikk-symbola. Det desse bevisa viser er \emph{ikkje} at Formlæra er sann (ei notasjon kan ikkje prove sanning om verda), men at ho er \emph{internt konsistent}. Kvart teorem følgjer frå tidlegare teorem og definisjonar utan motseiing.

Dette er eit sterkt resultat. Ein filosof som les desse bevisa kan gjera tre ting: kontrollere kvart skritt, foreslå alternative bevis, eller identifisere skritt som treng fleire premissar. Formskrifta gjer desse kritiske operasjonane \emph{mogelege}. Utan ho ville lesaren måtte rekonstruere bevisgangen frå prosa, og kvar slik rekonstruksjon ville variere.

Den djupare systemteoretiske innsikta er denne: \emph{ein teori som kan uttrykast formelt, kan også falsifiserast formelt}. Bevisgangane ovanfor kan sjekkast mekanisk. Formlæra har flytta seg frå ein filosofisk posisjon til eit falsifiserbart system. Det er den avgjerande progresjonen ein formell notasjon leverer.

\end{document}
